% Generated draft: requires mathematical and source review.
\documentclass[11pt]{article}
\usepackage[margin=25mm]{geometry}
\usepackage{fontspec}
\setmainfont{Latin Modern Roman}
\usepackage{amsmath,amssymb,mathtools,unicode-math}
\setmathfont{Latin Modern Math}
\usepackage{enumitem,needspace,xurl,hyperref}
\setlength{\parindent}{0pt}
\setlength{\parskip}{5pt}
\setlength{\emergencystretch}{3em}
\newcommand{\N}{\mathbb{N}}
\newcommand{\Z}{\mathbb{Z}}
\newcommand{\Q}{\mathbb{Q}}
\newcommand{\R}{\mathbb{R}}
\newcommand{\C}{\mathbb{C}}
\newcommand{\F}{\mathbb{F}}
\newcommand{\A}{\mathbb{A}}
\newcommand{\PP}{\mathbb{P}}
\newcommand{\E}{\mathbb{E}}
\newcommand{\eps}{\varepsilon}
\newcommand{\dd}{\,\mathrm{d}}
\begin{document}
\section*{UnsolvedMath ID: 30002434}
\textbf{Source problem code:} catalogue entry
\subsection{Definitions and mathematical statement}
Under the conventional interpretation of the notation, let \(r\in\Z\) with \(r\geq3\), and let \(h\in\Z\setminus\{0\}\) be fixed independently of \(X\). Write \(\Z_{>0}=\{1,2,\ldots\}\). For \(m\in\Z_{>0}\), define the \(r\)-fold divisor function by
\[
d_r(m)=\#\bigl\{(a_1,\ldots,a_r)\in\Z_{>0}^{\,r}:a_1\cdots a_r=m\bigr\}.
\]
For real \(X\geq1\), interpret the shifted correlation as
\[
D_r(X;h)=\sum_{\substack{n\in\Z_{>0}\\ n\leq X,\ n+h\geq1}}d_r(n)d_r(n+h).
\]
The supplied record does not explicitly define \(d_r\), specify the integer domains of \(r,h,n\), or state a convention when \(n+h\leq0\). The definitions above make these conventional choices explicit; their agreement with the intended source formulation requires review.

\textbf{Conjecture as supplied, with these conventions.} For every integer \(r\geq3\) and every fixed nonzero integer \(h\), there exists a polynomial \(P_{r,h}\in\R[u]\), denoted \(P_{2(r-1)}(u;h)\) in the supplied record, of degree exactly \(2(r-1)\), such that
\[
D_r(X;h)\sim X P_{r,h}(\log X)\qquad(X\to+\infty).
\]
Here \(\log\) is the natural logarithm, the polynomial may depend on both \(r\) and \(h\) but not on \(X\), and \(\sim\) means
\[
\lim_{X\to+\infty}\frac{D_r(X;h)}{X P_{r,h}(\log X)}=1.
\]
No uniformity as \(r\) or \(h\) varies is specified. This ratio asymptotic does not determine the lower-degree coefficients of \(P_{r,h}\). If a full polynomial expansion with an error smaller than \(X\) is intended, that stronger assertion is not specified by the displayed source statement and must be clarified.
\subsection{Short English statement}
For each fixed integer \(r\geq3\) and fixed nonzero integer shift \(h\), the supplied conjecture asks whether the correlation of the \(r\)-fold divisor function with its shift is asymptotic to \(X\) times a polynomial in \(\log X\) of degree \(2(r-1)\). Integer domains and the treatment of nonpositive shifted arguments are explicit conventional interpretations requiring review. The displayed ratio asymptotic leaves lower polynomial coefficients undetermined.
\subsection{Sources}
\begin{itemize}
\item Local registry: \texttt{data/unsolvedmath/problems.json}, supplied record with permanent UnsolvedMath ID \(30002434\), source problem code \texttt{catalogue entry}, title \emph{Polynomial Asymptotics for Shifted Divisor Correlations}.
\item Supplied report citation: \emph{Analytic Number Theory (2014)}, Oberwolfach Reports, DOI \texttt{10.4171/owr/2013/51}, pp.~2963--3037. Supplied source URL and literature-source link labeled ``original OWR report'': \url{https://doi.org/10.4171/owr/2013/51}. The background attributes the record to Oberwolfach workshop open-problem discussions. No individual proposer is supplied. The proposed year and year tag are 2013, whereas the supplied report citation says 2014; this dating discrepancy is unresolved here.
\item Supplied literature-source link labeled ``current additive-divisor literature'': \url{https://ems.press/content/serial-article-files/53325}. No author, article title, or precise supporting location is supplied for this link.
\item Literature assessment dated exactly \texttt{2026-08-21}: the supplied curation labels the problem open, describes the \(r=2\) case as classical, and states that the conjectured polynomial asymptotic for \(r\geq3\) remains open despite function-field analogues and later partial progress. This is supplied dated curation derived from the Oberwolfach record, not an independently verified literature assessment; its time-sensitive status and cited claims require independent checking.
\item The registry labels the statement \texttt{corrected\_verified}, with the supplied explanation that mathematical content was unchanged and notation, typography, and wording were normalized. This records the registry's characterization and does not establish independent mathematical or source verification of this draft.
\end{itemize}
\end{document}
